\scanpage{11}
\mspage{5}
For the amalgamated sum to exist in \deleted{$\operatorname{Esp}$} $\mathcal E$ and for $X\mathrel{\mapstochar\relbar}|X|$ to commute with it, it is necessary and sufficient that there exist a finite family of \inserted{pairs of} functions
\[
\begin{array}{l}
f_\alpha:\deleted{g}\ X_1\mathbin{-}R\\
g_\alpha:X_2\mathbin{-}R
\end{array}
\]
such that one has
\[
\text{a)}\quad f_\alpha(i_1(z))=g_\alpha(i_2(z))\qquad\forall\alpha\in A,\ z\in Z
\]
b) \deleted{If $x,y$ the $f_\alpha$ \unsure{x} $(f_\alpha):$} the $f_\alpha$ ($g_\alpha$) separate the points of $X_1$ ($X_2$) i.e.
\[
 X_1\mathbin{-}R^I\qquad(X_2\mathbin{-}R^A)
\]
are injective
\[
\begin{array}{ll}
\text{c)}&\text{If }x_1\deleted{\unsure{\alpha}}\in X_1\smallsetminus i_1(Z)\\
&x_2\in X_2\smallsetminus i_2(Z)\\
&\exists\alpha\text{ with }f_\alpha(x_1)\ne g_\alpha(x_2)
\end{array}
\]

\uline{NB} One can generalize, under the conditions of the proposition, there exists the tame structure one wants on $X$, if and only if $\exists$ a finite family $(F_\alpha)$ \deleted{of} of inductive systems
\[
 (f_{\alpha i})_{i\in I}\qquad f_{\alpha i}:X_i\underset{\text{tame}}{\longrightarrow}R),
\]
obtaining
\[
 f_\alpha:X\longrightarrow R
\]
such that the $f_\alpha$ separate the points of $X$, i.e.
\[
 f=(f_\alpha):X\hookrightarrow R^A
\]
\mspage{10}
